% ---------------------------------------------------------------------------
%  Motion cues for rotation frames.
%
%  Beamer overlays are jump-cuts by nature -- one frame replaces the last.
%  These macros fake continuous motion across that cut with two tricks:
%    1. GHOSTING: redraw the pivot's PRE-rotation position at low opacity
%       behind the post-rotation tree, connected by a faint dashed path.
%       Flipping between the two overlay steps then reads as the ghost
%       "solidifying" into its new spot, instead of a hard swap.
%    2. PIVOT ARCS: a curved arrow swept around the pivot node showing the
%       rotation's direction, so which way "up" happens is never ambiguous.
% ---------------------------------------------------------------------------

% \rotarc{pivot coord}{radius}{start angle}{end angle}
% Curved arrow around a pivot, drawn in the Fibonacci-gold rotation colour.
\newcommand{\rotarc}[4]{%
  \draw[avlgoldarrow, line width=1.2pt]
    (#1) ++(#3:#2) arc (#3:#4:#2);%
}
% Same, in the alert colour -- for a deletion-triggered rotation.
\newcommand{\rotarcalert}[4]{%
  \draw[draw=avlalert, -{Stealth[length=4pt]}, line width=1.2pt]
    (#1) ++(#3:#2) arc (#3:#4:#2);%
}

% \pivotdot{coord} -- small solid marker pinning down the exact pivot point.
\newcommand{\pivotdot}[1]{%
  \fill[avlfib] (#1) circle (1.6pt);%
}

% \ghostnode{x}{y}{label} -- a faded afterimage of a node at its OLD spot.
\newcommand{\ghostnode}[3]{%
  \node[circle, draw=avlmuted, fill=avlground, text=avlmuted, opacity=0.32,
        inner sep=1pt, minimum size=6.5mm, font=\footnotesize] at (#1,#2) {#3};%
}
% \ghostpath{from}{to} -- faint dashed thread linking an old spot to the new one.
\newcommand{\ghostpath}[2]{%
  \draw[avlmuted, opacity=0.35, dashed, line width=0.5pt] (#1) -- (#2);%
}
